\begin{pastexam}[label={px:ncku-112-prog-12}]{成大資訊聯招 112 學年度 第 12 題}
（填空）

(10\%) The following procedure is a bottom-up method for matrix chain order. This procedure assumes that
matrix $A_i$ has dimensions $p_{i-1}\times p_i$ for $i=1,2,\ldots,n$. Its input is a sequence
$p=\langle p_0,p_1,\ldots,p_n\rangle$, where $p.length=n+1$. Let $m[i,j]$ be the minimum number of scalar
multiplications needed to compute the matrix $A_iA_{i+1}\cdots A_j$. The procedure uses an auxiliary table
$m[1..n,1..n]$ for storing the $m[i,j]$ costs and another auxiliary table $s[1..n-1,2..n]$ that records
which index of $k$ achieved the optimal cost in computing $m[i,j]$. Please fill in the empty statements.

\begin{quote}
$n=p.length-1$\\
let $m[1..n,1..n]$ and $s[1..n-1,2..n]$ be new tables\\
\textbf{For} $i=1$ \textbf{to} $n$\\
\hspace*{2em}$m[i,i]=0$\\
\textbf{For} $l=2$ \textbf{to} $n$ \quad // $l$ is the chain length\\
\hspace*{2em}\textbf{For} $i=1$ \textbf{to} $n-l+1$\\
\hspace*{4em}$j=$ \blank{a}~(3\%)\\
\hspace*{4em}$m[i,j]=\infty$\\
\hspace*{4em}\textbf{For} $k=i$ \textbf{to} $j-1$\\
\hspace*{6em}$q=$ \blank{b}~(4\%)\\
\hspace*{6em}\textbf{If} $q<m[i,j]$\\
\hspace*{8em}$m[i,j]=q$\\
\hspace*{8em}$s[i,j]=$ \blank{c}~(3\%)
\end{quote}
\end{pastexam}

\begin{pastsoln}
答案：(a) $i + l - 1$；(b) $m[i, k] + m[k+1, j] + p_{i-1}\,p_k\,p_j$；(c) $k$。

這是資結書第~18.2~節的矩陣鏈乘（CLRS 的 \textsc{Matrix-Chain-Order}）。

\textbf{(a)}。 $l$ 是鏈的長度：$A_i \cdots A_j$ 有 $j - i + 1 = l$ 個矩陣，所以 $j = i + l - 1$。 外層依 $l$ 由小到大，
算 $m[i, j]$ 時用到的 $m[i, k]$、$m[k+1, j]$ 都是更短的鏈，已經算好（由下而上的順序）。 $i$ 跑到 $n - l + 1$ 時
$j = n$，剛好不越界。

\textbf{(b)}。 在 $A_k$ 和 $A_{k+1}$ 之間切開：先算 $A_i \cdots A_k$（$p_{i-1} \times p_k$，最少 $m[i, k]$ 次）和
$A_{k+1} \cdots A_j$（$p_k \times p_j$，最少 $m[k+1, j]$ 次），兩個結果相乘要 $p_{i-1} p_k p_j$ 次純量乘法：
\[
q = m[i, k] + m[k + 1, j] + p_{i-1}\, p_k\, p_j.
\]
最佳解一定在某個 $k$ 切開，兩段又各自必須最佳（最佳子結構），所以
$m[i, j] = \min_{i \le k < j} q$。

\textbf{(c)}。 $s[i, j]$ 記下達到最小值的切點，填 $k$。 之後由 $s[1, n]$ 遞迴往下就能印出最佳括號方式。

小例子核對：$p = \langle 10, 100, 5, 50\rangle$。 $m[1,2] = 10 \cdot 100 \cdot 5 = 5000$，$m[2,3] = 100 \cdot 5 \cdot 50 = 25000$。
$m[1,3]$：$k = 1$ 時 $0 + 25000 + 10 \cdot 100 \cdot 50 = 75000$；$k = 2$ 時 $5000 + 0 + 10 \cdot 5 \cdot 50 = 7500$。
取 $7500$，$s[1, 3] = 2$，也就是 $((A_1A_2)A_3)$。
\end{pastsoln}
